\begin{frame}[t]{Giải thích thước đo dẫn nguồn: Validity và F1}
\purpose{Đánh giá dẫn nguồn theo hai tầng độc lập: tính hợp lệ cú pháp và độ chính xác đối chiếu với bằng chứng chuẩn.}
\begin{columns}[T,onlytextwidth]
\begin{column}{.48\textwidth}
\heading{1. Citation Validity (Hợp lệ cú pháp)}
{\fontsize{13.5}{17.5}\selectfont
Kiểm tra các chỉ số $[i]$ trong câu trả lời có thuộc phạm vi ngữ cảnh được cấp $\{1, \dots, k\}$ hay không:
\vspace{-2mm}
\[
\text{Validity} = \frac{|\mathcal{C}_{\text{hợp lệ}}|}{|\mathcal{C}_{\text{hợp lệ}}| + |\mathcal{C}_{\text{không hợp lệ}}|}
\]
\vspace{-2mm}
\begin{itemize}
  \setlength{\itemsep}{2pt}
  \item \textbf{Mục đích:} Phát hiện ảo giác cú pháp (chỉ số ngoài dải, ví dụ dẫn $[8]$ khi chỉ cấp $k=5$ đoạn).
  \item \textbf{Ý nghĩa:} Validity $= 1{,}0$ chỉ bảo đảm chỉ số tồn tại, \textit{chưa bảo đảm} dẫn đúng sự thật.
\end{itemize}\par}
\end{column}
\begin{column}{.48\textwidth}
\heading{2. Citation F1 (Khớp bằng chứng chuẩn)}
{\fontsize{13.5}{17.5}\selectfont
Ánh xạ các chỉ số hợp lệ sang tập Chunk ID $\mathcal{C}$ và đối chiếu với tập chunk chuẩn $\mathcal{G}$ (gold evidence):
\vspace{-2mm}
\[
P_{\text{cit}} = \frac{|\mathcal{C} \cap \mathcal{G}|}{|\mathcal{C}|},\qquad R_{\text{cit}} = \frac{|\mathcal{C} \cap \mathcal{G}|}{|\mathcal{G}|}
\]
\vspace{-4mm}
\[
\text{Citation F1} = \frac{2 \cdot P_{\text{cit}} \cdot R_{\text{cit}}}{P_{\text{cit}} + R_{\text{cit}}}
\]
\vspace{-2mm}
\begin{itemize}
  \setlength{\itemsep}{2pt}
  \item \textbf{Mục đích:} Đảm bảo câu trả lời neo đúng bằng chứng thực tế, không dẫn đoạn gây nhiễu (distractor).
  \item \textbf{Đặc điểm:} Tính tất định theo tập hợp; bằng 0 nếu dẫn nhầm đoạn dù cú pháp hợp lệ.
\end{itemize}\par}
\end{column}
\end{columns}
\vspace{2mm}
\takeaway{\textbf{Validity} kiểm tra trích dẫn có tồn tại trong context; \textbf{Citation F1} đo lường trích dẫn có trỏ đúng bằng chứng chuẩn.}
\slidenote{Cả hai thước đo đều được tính tất định (deterministic), không phụ thuộc vào LLM judge.}
\end{frame}
